\chapter{Rhabdomyolysis}
\index{Rhabdomyolysis}
\label{sec:Rhabdomyolysis}

\section{Overview}
Rhabdomyolysis is a syndrome of acute skeletal muscle tissue death with release of intracellular muscle contents (myoglobin, CK, potassium, phosphate, uric acid) into the circulation.
\section{Causes}
Causes include crush injuries, prolonged immobilization, extreme exertion, seizures, cancerous hyperthermia, medications (statins, fibrates, antipsychotics), electrolyte imbalances, infections, and inflammatory myopathies. The underlying mechanisms involve complex interactions between genetic predisposition and environmental factors that disrupt normal physiological processes. The underlying mechanisms involve complex interactions between genetic predisposition and environmental factors. Disruption of normal physiological processes leads to the characteristic features of the condition. Multiple contributing factors may act synergistically to trigger or worsen the condition. Understanding the specific cause helps guide targeted treatment approaches.
\section{Risk Factors}


\begin{itemize}[noitemsep]
  \item Genetic predisposition and family history
  \item Advancing age
  \item Lifestyle factors (diet, exercise, smoking, alcohol)
  \item Environmental exposures
  \item Underlying medical conditions
  \item Medication use
\end{itemize}
\section{Signs and Symptoms}

\subsection{Common symptoms}
\begin{itemize}[noitemsep]
  \item Symptoms vary depending on the severity and extent of the condition Pain or discomfort in the affected area Functional impairment affecting daily activities Changes in appearance or function of affected tissues Systemic symptoms may include fatigue, fever, or weight changes Symptoms may develop gradually or appear suddenly
  \item Pain or discomfort in the affected area
  \end{itemize}

\subsection{Emergency warning signs}
\begin{itemize}[noitemsep]
  \item Severe or worsening symptoms of rhabdomyolysis require immediate medical evaluation
\end{itemize}
\section{Progression and Stages}
The condition may progress through different stages of severity over time. Early stages often involve mild or subtle symptoms that may be overlooked. Moderate stages involve more noticeable symptoms affecting daily function. Advanced stages may involve significant impairment and complications.

\begin{itemize}[noitemsep]
  \item Your outlook depends on severity and promptness of treatment
  \item Most patients recover with adequate hydration.
\end{itemize}
\section{Complications}
\begin{itemize}[noitemsep]
  \item You may experience muscle pain, muscle weakness, dark (cola-colored) urine, and swelling of affected muscles
  \item Complications include acute kidney injury (myoglobin nephrotoxicity), hyperkalemia (irregular heartbeats), disseminated intravascular coagulation, and compartment syndrome.
  \item Disease progression leading to worsening symptoms
  \item Secondary infections or injuries
  \item Side effects of treatment
  \item Reduced quality of life
  \item Emotional and psychological impact
\end{itemize}
\section{Diagnosis}
Doctors diagnose this based on elevated CK (typically >5 times the upper limit of normal, often >10,000 IU/L), elevated myoglobin, and elevated transaminases. Comprehensive medical history and physical examination Laboratory tests to assess organ function and rule out other conditions Imaging studies to visualize affected structures Specialized testing depending on the affected system Consultation with relevant specialists Monitoring of disease progression over time

\begin{itemize}[noitemsep]
  \item Comprehensive medical history and physical examination
  \item Laboratory tests to assess organ function
  \item Imaging studies to visualize affected structures
  \item Specialized testing depending on the affected system
  \item Consultation with relevant specialists
\end{itemize}
\section{Prevention}
\begin{itemize}[noitemsep]
  \item Treatment focuses on aggressive IV fluid resuscitation (target urine output 200--300 mL/h), urinary alkalinization (to prevent myoglobin precipitation), correction of electrolyte abnormalities, and monitoring for complications
  \item Kidney replacement therapy may be necessary for severe AKI.
  \item Maintain a healthy lifestyle with balanced diet and exercise
  \item Avoid known risk factors
  \item Attend regular medical check-ups for early detection
  \item Follow recommended screening guidelines
\end{itemize}
\section{Treatment Options}
Medications to manage symptoms and address underlying mechanisms Lifestyle modifications and self-care measures Physical therapy and rehabilitation Surgical intervention when indicated Regular monitoring and follow-up care Patient education and support services

\begin{itemize}[noitemsep]
  \item Medications to manage symptoms and address underlying mechanisms
  \item Lifestyle modifications and self-care measures
  \item Physical therapy and rehabilitation
  \item Surgical intervention when indicated
  \item Regular monitoring and follow-up care
\end{itemize}
\section{Living with It}
\begin{itemize}[noitemsep]
  \item Follow your treatment plan as prescribed
  \item Maintain regular follow-up with your healthcare provider
  \item Make necessary lifestyle adjustments
  \item Seek support from family, friends, or support groups
  \item Stay informed about your condition
\end{itemize}
\section{Outlook}
The outlook varies depending on the specific condition, severity at diagnosis, and response to treatment. Early intervention generally leads to better outcomes. Many conditions can be effectively managed with appropriate treatment. Regular follow-up care is important for monitoring and treatment adjustment.
